\section{可插拔后端：\texttt{vma\_ops\_t} 接口}

与 PMM（第~\ref{ch:pmm}~章）和 kmalloc（第~\ref{ch:kmalloc}~章）完全相同的
"dispatch + 后端"模式：

\codefile{src/include/kernel/vma.h}
\begin{lstlisting}[style=cstyle]
typedef struct vma_ops {
    const char* name;
    void (*init)(void);
    int  (*add)(uint32_t start, uint32_t end,
                uint32_t flags, const char* name);
    int  (*remove)(uint32_t start);
    const vm_area_t* (*find)(uint32_t addr);
    void (*dump)(void);
    unsigned (*count)(void);
} vma_ops_t;
\end{lstlisting}

\begin{figure}[H]
  \centering
  % figures/ch11/fig02-vma-dispatch.tex — auto-extracted from chapter source
\begin{tikzpicture}[
  box/.style={draw, rounded corners=3pt, minimum height=0.7cm,
              minimum width=3.5cm, font=\small, align=center},
  arr/.style={-{Stealth[length=5pt]}, thick},
]
  \node[box, fill=blue!8]   (caller) at (0, 0)
    {调用者 (kmain / vmm / shell)};
  \node[box, fill=yellow!12](api)    at (0, -1.3)
    {\texttt{vma\_add()} / \texttt{vma\_find()}};
  \node[box, fill=orange!12](disp)   at (0, -2.6)
    {\texttt{vma.c} dispatch};
  \node[box, fill=green!12] (sa)     at (-4.0,-4.2)
    {\texttt{vma\_sorted\_array.c}\\（排序数组）};
  \node[box, fill=green!12] (rb)     at (0,-4.2)
    {\texttt{vma\_rbtree.c}\\（红黑树）};
  \node[box, fill=green!12] (mt)     at (4.0,-4.2)
    {\texttt{vma\_maple.c}\\（Maple Tree）};
  \draw[arr] (caller)--(api);
  \draw[arr] (api)--(disp);
  \draw[arr] (disp)--(sa);
  \draw[arr] (disp)--(rb);
  \draw[arr] (disp)--(mt);
\end{tikzpicture}

  \caption{\texttt{vma\_ops\_t} 可插拔后端调用链——三种后端通过 \texttt{kconfig.h} 选择}
  \label{fig:vma-dispatch}
\end{figure}

% ============================================================================

